\subsection{\texorpdfstring{群的 \(\tau\)-扩张}{群的 \textbackslash tau-扩张}}

在本小节中，我们固定一个群 G、一个以乘法记写的交换群 F，以及一个群同态 \(\tau\)（从 G 到 F 的自同构群 \(\operatorname{Aut}(F)\) 的）。我们把 G 的单位元素记作 e，把 F 的单位元素记作 1。

回忆（I, §6, No.~1, 第 65 页）一个扩张 \(\mathcal{E}\)（\(G\) 被 \(F\) 的）是指一个三元组 \((\Gamma, \iota, \pi)\)，其中 \(\Gamma\) 是一个群，\(\pi: \Gamma \to G\) 是一个满射同态，而 \(\iota\) 是一个从 \(F\) 到 \(\Gamma\) 的单射同态，使得 \(\operatorname{Im}(\iota) = \operatorname{Ker}(\pi)\)。设 \(\mathcal{E} = (\Gamma, \iota, \pi)\) 是这样一个扩张。对每个 \(\gamma \in \Gamma\)，映射 \(\varphi_{\gamma}: F \to F\)（由

\[
\iota (\varphi_ {\gamma} (f)) = \gamma \iota (f) \gamma^ {- 1}
\]

对 \(f \in \mathrm{F}\) 所定义的）是 \(\mathrm{F}\) 的一个自同构。由于 \(\mathrm{F}\) 是交换的，对每个 \(f \in \mathrm{F}\)，由 \(\iota(f)\) 所定义的自同构是 \(\mathrm{F}\) 上的恒等映射。通过过渡到商，我们得到一个同态 \(\operatorname{Int}_{\mathcal{E}}\)（从 \(\mathrm{G}\) 到 \(\operatorname{Aut}(\mathrm{F})\) 的），它由

\[
\gamma \iota (f) \gamma^ {- 1} = \iota (\mathrm{Int} _ {\mathcal {E}} (\pi (\gamma)) \cdot f)
\]

（对 \(\gamma \in \Gamma\) 与 \(f\in \mathrm{F}\)）刻画。

G 被 F 的一个 \(\tau\)-扩张是指一个扩张 \(\mathcal{E} = (\Gamma, \iota, \pi)\)，它使得 \(Int_{\mathcal{E}}\) 等于 \(\tau\)，换句话说，它满足关系

\[
\gamma \iota (f) \gamma^ {- 1} = \iota (\tau (\pi (\gamma)) \cdot f)\tag{1}
\]

（对 \(\gamma \in \Gamma\) 与 \(f \in \mathrm{F}\)）。若 \(\mathcal{E} = (\Gamma, \iota, \pi)\) 与 \(\mathcal{E}' = (\Gamma', \iota', \pi')\) 是 \(\tau\)-扩张，则一个 \(\tau\)-扩张的态射（从 \(\mathcal{E}\) 到 \(\mathcal{E}'\) 的）是指一个从 \(\mathcal{E}\) 到 \(\mathcal{E}'\) 的扩张的态射（I, §6, No.~1, 第 65 页），即一个群同态 \(u: \Gamma \to \Gamma'\)，使得 \(\pi' \circ u = \pi\) 且 \(\iota' = u \circ \iota\)。注意 \(\tau\)-扩张构成一个结构种，其每个态射都是同构（集合论, IV, §1, No.~5, 第 264 页与 I, §6, No.~1, 第 66 页，命题 1）。用 \(\operatorname{Iso}_{\tau} \{\mathcal{E}, \mathcal{F}\}\) 表示关系

“\(\mathcal{E}\) 与 \(\mathcal{F}\) 是同构的 \(\tau\)-扩张。”

这是一个等价关系；扩张 \(\mathcal{E}\) 的类是与 \(\mathcal{E}\) 等价（对 \(\mathrm{Iso}_{\tau}\) 而言）的对象的类（集合论, II, §6, No.~6, 第 122 页）。

引理 1. —— 关系

\textbf{\texorpdfstring{“α 是对 \(Iso_{\tau}\) 而言一个 τ-扩张的类”}{“α 是对 Iso\_\{\textbackslash tau\} 而言一个 τ-扩张的类”}}

在 \(\alpha\) 中是集合化的。

设 \(E_{0}=F\times G\)，并考虑映射 \(\iota_{0}:F\to E_{0}\) 与 \(\pi_{0}:E_{0}\to G\)，它们由 \(\iota_{0}(f)=(f,e)\)（对 \(f\in F\)）与 \(\pi_{0}(f,g)=g\)（对 \((f,g)\in F\times G\)）定义。设 \(\mathcal{E}=(\Gamma,\iota,\pi)\) 是 G 被 F 的一个 \(\tau\)-扩张。映射 \(\pi\) 是满射的，故有一个截面 \(\sigma:G\to\Gamma\)，使得 \(\sigma(e)=e\)。映射 \(u:(f,g)\mapsto\iota(f)\sigma(g)\)（从 \(F\times G\) 到 \(\Gamma\) 的）是双射的。我们给 \(F\times G\) 装备由结构转移得到的群律。于是三元组 \((\mathrm{E}_{0},\iota_{0},\pi_{0})\) 是一个同构于 E 的 \(\tau\)-扩张。引理 1 于是由集合论, II, §6, No.~6, 第 122 页得出。

我们用 \(\operatorname{Ex}_{\tau}(\mathbf{G},\mathbf{F})\) 表示 \(\tau\)-扩张的类（对关系 \(Iso_{\tau}\) 而言的）所成之集。

例. —— 1) 外半直积 \((\mathrm{F} \times_{\tau} \mathrm{G}, i, p)\)（I, §6, No.~1, 第 66 页，命题 3）是一个 \(\tau\)-扩张；我们把它记作 \(\mathcal{I}_{\tau}\)。我们说一个 \(\tau\)-扩张是半平凡的，如果它同构于 \(\tau\)-扩张 \(\mathcal{I}_{\tau}\)。

\begin{enumerate}
\def\labelenumi{\arabic{enumi})}
\setcounter{enumi}{1}
\tightlist
\item
  对 \(i \in \{1,2\}\)，取一个群 \(G_{i}\)、一个交换群 \(F_{i}\)，以及一个群同态 \(\tau_{i}\)（从 \(G_{i}\) 到 \(F_{i}\) 的自同构群的）。我们用 \(\tau_{1} \times \tau_{2}: G_{1} \times G_{2} \to \text{Aut}(F_{1} \times F_{2})\) 表示由
\end{enumerate}

\[
(\tau_ {1} \times \tau_ {2}) (g _ {1}, g _ {2}) \cdot (f _ {1}, f _ {2}) = (\tau_ {1} (g _ {1}) \cdot f _ {1}, \tau_ {2} (g _ {2}) \cdot f _ {2})
\]

（对 \(g_{1} \in \mathrm{G}_{1}\)、\(g_{2} \in \mathrm{G}_{2}\)、\(f_{1} \in \mathrm{F}_{1}\) 与 \(f_{2} \in \mathrm{F}_{2}\)）所定义的同态。设 \(\mathcal{E}_1 = (\Gamma_1, \iota_1, \pi_1)\) 是一个 \(\tau_1\)-扩张（\(\mathrm{G}_1\) 被 \(\mathrm{F}_1\) 的），\(\mathcal{E}_2 = (\Gamma_2, \iota_2, \pi_2)\) 是一个 \(\tau_2\)-扩张（\(\mathrm{G}_2\) 被 \(\mathrm{F}_2\) 的）。则三元组 \((\Gamma_1 \times \Gamma_2, \iota_1 \times \iota_2, \pi_1 \times \pi_2)\) 是一个 \(\tau_1 \times \tau_2\)-扩张（\(\mathrm{G}_1 \times \mathrm{G}_2\) 被 \(\mathrm{F}_1 \times \mathrm{F}_2\) 的）；它称为扩张 \(\mathcal{E}_1\) 与 \(\mathcal{E}_2\) 的外积，记作 \(\mathcal{E}_1 \times \mathcal{E}_2\)。

